\section{An explicit exponent for broad Gaussian windows}
\label{sec:broad}
Write the centered Gaussian-smoothed limiting vector in the form
\[
   Y_h(z)=2\operatorname{Re}\sum_{\gamma>0}
                   e^{-h\gamma^2} b_\gamma z_\gamma ,
       \qquad h>0,
\]
where $z$ is Haar distributed on the closure of the orbit
$(e^{i\gamma t})_{\gamma>0}$. The vectors $b_\gamma\in\mathbb C^d$
are the unsmoothed coefficients, with coincident ordinates grouped.
Thus $Y_h=F_z(0)-b_0$ in the notation of the prime-power formula.
A common time translation leaves the law unchanged. The sum of
$\|b_\gamma\|$ over $0<\gamma\le T$ is $O_q((1+T)\log(2+T))$ for
the vectors $b_\gamma(a)=-d^{-1}\sum_\chi
\overline{\chi(a)}m_\chi(\gamma)$.

In increasing order of $\gamma$, retain $\gamma$ if adjoining
$\operatorname{Re}b_\gamma$ and $\operatorname{Im}b_\gamma$ increases
the dimension of the real span. Write $\lambda_1<\cdots<\lambda_r$
for the retained ordinates. The zero-divisor independence lemma implies
that this span eventually equals $\mathbb R^d$, so $1\le r\le d$.
For any nonzero real vector $v$ put
\[
 A_h(v)=\max_{1\le j\le r}e^{-h\lambda_j^2}
                         |v\cdot b_{\lambda_j}|>0 .
\]
Here the dot product extends the real dot product complex-linearly in
the second argument.

\begin{theorem}\label{thm:broad}
There are $h_0$ and $C$, depending only on this spectrum, such that
for $h\ge h_0$, every nonzero $v$, and every Borel $E\subset\mathbb R$,
\[
  \Pr\{v\cdot Y_h\in E\}
    \le \min\left\{1,\;C\left(\frac{|E|}{A_h(v)}\right)^{1/(2r)}\right\}.
\]
The characteristic function in this direction satisfies
\[
 \left|\mathbb E\exp(iu\,v\cdot Y_h)\right|
       \le C(1+|u|A_h(v))^{-1/(2r)} .
\]
In particular $1/(2d)$ is always an admissible exponent.
\end{theorem}

The directional normalization is necessary for a bound uniform in $h$,
since every centered projection tends to zero as $h\to\infty$.

\begin{proof}
Along any orbit translate, divide the scalar projection by $A_h(v)$.
Its retained part is a trigonometric polynomial at the fixed $2r$
frequencies $\nu\in\{\pm\lambda_1,\ldots,\pm\lambda_r\}$.
At least one of its complex coefficients has modulus one and all
have modulus at most one. At every time, the vector of its derivatives
of orders $1,\ldots,2r$ is the matrix $((i\nu)^k)_{k,\nu}$ times
the vector of its phased coefficients. Factoring $i\nu$ from each column leaves an invertible Vandermonde matrix.
The maximum of these derivative moduli is therefore bounded below
by a positive constant depending only on the retained ordinates.

For an omitted $\gamma$, the two real vectors
$\operatorname{Re}b_\gamma$ and $\operatorname{Im}b_\gamma$ belong
to the span at the preceding retained stage. Choose once and for
all a basis from the retained vectors at each of the finitely many
stages. Boundedness of each coordinate map gives
\[
 \frac{e^{-h\gamma^2}|v\cdot b_\gamma|}{A_h(v)}
 \le C_q\|b_\gamma\|\!
        \sum_{j:\lambda_j<\gamma}
                    e^{-h(\gamma^2-\lambda_j^2)} .
\]
For $h\ge1$, the right side times $(1+\gamma)^{2r+1}$ has a summable
majorant. Every term tends to zero, so the omitted part tends uniformly
to zero with all derivatives through order $2r+1$, uniformly in
$v$ and in the phases. For $h$ sufficiently large the normalized
function thus has some derivative of order $1,\ldots,2r$ bounded
away from zero at every point, and all derivatives through order
$2r+1$ are bounded above uniformly. The one-dimensional finite-type
estimates, applied on unit orbit segments and averaged over Haar
measure, give the asserted bounds.
\end{proof}

For each fixed $h\ge h_0$, the late unit-interval estimates for the
arithmetic signal have the same exponent after division by $A_h(v)$.
To see this, normalize $\|v\|=1$. The map $v\mapsto A_h(v)$ is continuous
and positive on the unit sphere, so its minimum is positive for this
fixed $h$. The differentiated explicit-formula remainder through order
$2r+1$ is consequently $o(1)$ even after division by $A_h(v)$, uniformly
over such $v$ as the left endpoint $T$ tends to infinity. The threshold
on $T$ in this assertion may depend on $h$.

\begin{corollary}[One projected oscillation at broad bandwidth]
For fixed $v\ne0$, let $\gamma_* >0$ be the first frequency such that
$v\cdot b_{\gamma_*}\ne0$. Then
\[
 \frac{v\cdot Y_h}{2e^{-h\gamma_*^2}|v\cdot b_{\gamma_*}|}
       \ \xrightarrow[h\to\infty]{\mathrm{law}}\ \cos\Theta,\qquad
       \Theta\ \text{uniform on }[0,2\pi].
\]
\end{corollary}

\begin{proof}
The projection of the phase group on this frequency is the
whole circle and the normalized tail converges uniformly to zero.
If $\gamma'>\gamma_*$ is the next contributing frequency, this coupling
has error $O_{q,v}(e^{-h((\gamma')^2-\gamma_*^2)})$ for $h\ge1$.
The limiting density $1/(\pi\sqrt{1-y^2})$ on $(-1,1)$ assigns mass
of order $\sqrt\varepsilon$ to $[1-\varepsilon,1]$. In the scalar
case the exponent $1/2$ is consequently sharp for bounds uniform
as $h\to\infty$.
\end{proof}

\begin{corollary}[Counting crossings of a smoothed prime statistic]
\label{cor:crossings}
Fix a nonzero real $v$. For all sufficiently large fixed $h$, the
number of crossings of its own limiting mean by either scalar arithmetic
statistic $v\cdot \mathcal P(t)$ or $v\cdot \mathcal Q(t)$ in the
log-time interval $[0,T]$ equals
\[
                    \frac{\gamma_*}{\pi}T+O_{q,v,h}(1).
\]
Every sufficiently late crossing is simple. If
$\phi=\arg(v\cdot b_{\gamma_*})$, there is exactly one crossing near
each sufficiently late point $(\pi/2+k\pi-\phi)/\gamma_*$.
\end{corollary}

\begin{proof}
Set $\theta=\gamma_*t+\phi$, subtract the limiting mean of the
chosen scalar statistic, and divide by
$2e^{-h\gamma_*^2}|v\cdot b_{\gamma_*}|$.
In this coordinate the first oscillation is $\cos\theta$. The sum of
all higher oscillations and its first derivative tend uniformly to
zero with $h$, by the preceding tail estimate with an extra weight
$1+\gamma/\gamma_*$. For fixed sufficiently large $h$, the additional
arithmetic remainder and its derivative tend to zero with $t$. The two
errors together can therefore be bounded in modulus, separately for
derivative and value, by $1/4$ for all sufficiently late times.
Outside the intervals $|\theta-(\pi/2+k\pi)|<\pi/6$ we have
$|\cos\theta|\ge1/2$. Inside each interval, $|\sin\theta|\ge\sqrt3/2$;
the perturbed cosine is strictly monotone there and has opposite signs
at the endpoints. Each such interval contributes exactly one simple
crossing. The arithmetic statistic is real analytic and nonconstant,
so a bounded initial interval contributes finitely many crossings.
\end{proof}

